\documentclass[10pt]{article}
\usepackage[margin=20mm]{geometry}
\usepackage{amsmath,amssymb,mathtools,fontspec,unicode-math,xurl,hyperref}
\IfFontExistsTF{DejaVu Serif}{\setmainfont{DejaVu Serif}}{\setmainfont{Latin Modern Roman}}
\setmathfont{latinmodern-math.otf}
\setlength{\parindent}{0pt}\setlength{\parskip}{6pt}\sloppy
\title{Technical verification of the geometry of permutation modules}
\author{AI-assisted technical review}
\date{}
\begin{document}
\maketitle
Abstract: This report examines four selected argument chains: the \(p\)-subgroup condition and additive quotient for modular fixed points, the normalizer step in point classification, the finite topological colimit, and the cyclic \(p\)-group boundary cases. The inspected steps retain their necessary normality, characteristic, finiteness, and prime-domain assumptions. Exact finite diagnostics agree with the corresponding elementary calculations. No substantive defect was established within this declared scope. This conclusion does not certify the full conservativity, stratification, or twisted-cohomology arguments, and is restricted to the acquired author manuscript. Scientific adjudication remains pending.\par
\section*{Identity and mathematical setting}
The reviewed work is Paul Balmer and Martin Gallauer, The geometry of permutation modules, arXiv:2307.04398v2, with a PDF version stamp of 21 July 2025, 76 pages [paper]. PDF page numbers agree with the printed numbers used here. The journal DOI associated with the acquisition record is 10.1007/s00222-025-01350-z. The journal Version of Record was not acquired, and textual equivalence to it is not assumed.\par
The domain is a finite group \(G\) and a field \(k\) of characteristic \(p>0\). The compact tensor triangulated category is the idempotent completion of the bounded homotopy category of finite permutation modules, denoted \(\mathcal K(G)\). Write \(G/\!/H=N_G(H)/H\) for the Weyl group. The modular fixed-point functor \(\Psi^H\), the ordinary fixed-vector functor, and the geometric localization have different meanings; replacing one by another would change the assertions being checked.\par
A same-conversation native helper supplied bounded proof analysis, static author-TeX excerpts, and independently written finite computations. The introduction and selected proof passages were read; full text extraction is not counted as full reading. Neither the author source nor author software was executed.\par
\section*{The \(p\)-subgroup condition and the additive quotient}
\subsection*{Necessity in Warning 5.1}
On p. 17 the manuscript explains why a tensor functor satisfying \[\Psi^H(kX)\simeq k(X^H)\] on finite \(G\)-sets cannot have arbitrary \(H\). Let \(P\) be a Sylow \(p\)-subgroup. The coaugmentation and augmentation \[k\longrightarrow k(G/P)\longrightarrow k\] have composite \([G:P]\), invertible in \(k\). Thus the tensor unit is a retract of \(k(G/P)\). A unital additive tensor functor to the stated nonzero target cannot send \(k(G/P)\) to zero. The fixed-set formula therefore forces \((G/P)^H\ne\varnothing\). An \(H\)-fixed coset \(gP\) gives \(g^{-1}Hg\le P\), so \(H\) is a \(p\)-group.\par
For the excluded extension, \(G=H=C_3\) over \(\mathbb F_2\) already exhibits the obstruction: \(k\) is a retract of \(kG\), but the regular \(G\)-set has no \(G\)-fixed point. This is a check of the necessity of the paper’s restriction, not a counterexample to its actual theorem.\par
\subsection*{The vanishing lemma retains normality}
Lemma 5.3 (p. 18) assumes a normal \(p\)-subgroup \(N\trianglelefteq G\), \(N\le H\), and \(N\nleq K\). After the stated double-coset reduction, a map \(k(G/H)\to k(G/K)\to k\) is determined by the transfer index \[[H:L],\qquad L=H\cap gKg^{-1}.\] The \(N\)-action on \(H/L\) has no fixed point. Indeed, a fixed coset would imply a conjugate of \(N\) lies in \(L\); normality in \(G\) then implies \(N\le K\), a contradiction. Every orbit consequently has size divisible by \(p\). Thus \(p\mid[H:L]\), and the map is zero over \(k\). Both normality and the characteristic hypothesis are used in the displayed proof.\par
\subsection*{Faithfulness and idempotent completion}
Proposition 5.4 on the same page identifies the permutation category for \(G/N\) with the additive quotient that kills modules \(k(G/K)\) with \(N\nleq K\). Inflation is fully faithful before quotienting. To check that the quotient introduces no new kernel between surviving objects, use rigidity to turn a factorization \(A\to B\) through the killed ideal into a factorization \(A\otimes B^\vee\to k\). Surviving transitive factors have stabilizers containing \(N\); a tensor product involving a killed stabilizer decomposes into factors still failing to contain \(N\). Lemma 5.3 kills each such factorization.\par
Every transitive module either survives as an inflated \(G/N\)-module or becomes zero. Full faithfulness lifts the remaining idempotents, which split in the completed source. This verifies the stated essential-surjectivity step. Equation (5.13) and Corollary 5.16 on p. 20 are consistent with this quotient and its inflation section. No defect was established in these steps. The general double-coset and Mackey formulas were used as mathematical dependencies, not re-audited from their original references.\par
\section*{Point classification and the normalizer step}
\subsection*{Conservativity-to-nilpotence reduction}
Proposition 7.1 (p. 26) considers a conservative coproduct-preserving tensor functor between big tensor triangulated categories. Rigidity reduces a morphism with compact source to one from the unit. Its tensor-power telescope maps to a sequential homotopy colimit of zero maps after applying the functor, and is therefore zero there. Conservativity makes the original telescope zero; compactness of the unit then forces vanishing at a finite stage. This verifies the specific telescope step. The full Conservativity Theorem 6.12 and the external spectrum-surjectivity theorem remain unproved dependencies of the classification.\par
\subsection*{Support direction and uniqueness}
The point \(\mathcal P_G(H,\mathfrak p)\) is the inverse image of a cohomological prime for \(G/\!/H\), as defined on p. 26. Lemma 5.21 (p. 22) gives the two alternatives behind Lemma 7.12 (pp. 27–28): \[\operatorname{kos}_G(K)\in\mathcal P_G(H,\mathfrak p)
 \quad\Longleftrightarrow\quad
 H\ \text{is subconjugate to }K.\] When \(H\) is not subconjugate to \(K\), the fixed-point image has the tensor unit as a direct summand and cannot belong to a proper prime. In the other case it is acyclic and becomes zero in the derived category used to define the point. Consequently \[\mathcal P_G(H,\mathfrak p)\subseteq\mathcal P_G(H',\mathfrak p')
 \quad\Longrightarrow\quad H'\ \text{is subconjugate to }H.\] Evaluating the test object with \(K=H\) checks this direction. Equality forces conjugacy, as Corollary 7.13 states.\par
For Proposition 7.14 (p. 28), set \(N=N_G(H)\). Corollary 4.9 (p. 16), invoked as a descent input, represents equal images by a prime over an intersection of conjugates of \(N\). Classification and naturality express that prime through a subgroup \(L\), with both \(L\) and its relevant conjugate \(N\)-conjugate to \(H\). Since \(H\trianglelefteq N\), these subgroups are exactly \(H\); the conjugating element must normalize \(H\), hence lies in \(N\). The comparison then occurs within \(N\), where inner conjugation acts trivially on the spectrum. Finally the section supplied by inflation in Corollary 5.16 gives equality of the original cohomological primes.\par
This checks the fusion step used for the uniqueness assertion of Theorem 7.16. The proof does not assume \(H\trianglelefteq G\), nor does it assert that unrestricted restriction maps or all modular fixed-point maps are injective. It uses the cohomological open and normality inside the normalizer. No missing normality assumption or reversed support inclusion was established in this chain.\par
\section*{The topology of the finite colimit}
Proposition 7.32 (p. 31) explicitly gives a set partition by locally closed strata. Recollection 10.7 (p. 41) then states the problem of attaching these strata topologically. The coproduct symbol in the partition is not a claim that the spectrum is a topological disjoint union.\par
Construction 11.1 and the maps in (11.5)–(11.8), pp. 41–42, use the category \(\mathcal E_p(G)\) of elementary abelian \(p\)-sections \((H,K)\), with \(K\trianglelefteq H\) and \(H/K\) elementary abelian. Quotient-type arrows are explicitly included. They are essential: the arrow \((H,L)\to(H,K)\), for \(K\le L\le H\), need not be an inclusion \(H/L\to H/K\).\par
The closed-map step in Theorem 11.10 (p. 43) is justified by finiteness. There are only finitely many objects \((H,K)\) because \(G\) is finite. If \(X\) is the topological colimit and \(C\subset X\) is closed, its inverse image \(C_a\) in each component is closed. Lemma 11.9 gives a closed component map \(\varphi_a\). Therefore \[\varphi(C)=\bigcup_{a\in\operatorname{Ob}\mathcal E_p(G)}\varphi_a(C_a)\] is a finite union of closed sets. The argument would not be valid for an arbitrary infinite family, but that is outside the theorem’s hypotheses.\par
The injectivity reduction also stays inside the indexing category. A point represented by \(\mathcal P_{H/K}(L/K,\mathfrak p)\) can be represented at \((H,L)\), since the quotient \(H/L\) is again elementary abelian. Equality of the global primes forces the kernel subgroups to be conjugate by Theorem 7.16. After identifying them, the invoked Quillen/descent relation in \(N_G(K)/K\) produces a common elementary abelian section of the form \((H\cap gH'g^{-1},K)\). Its conjugation/restriction arrows identify the two points already in the colimit. Together with the invoked surjectivity, this yields a continuous closed bijection and hence a homeomorphism.\par
The original Quillen and separable-descent theorems were not independently re-proved. The checked argument is the internal reduction, admissibility of its sections, and closed-map inference; it does not certify all later local ring presentations.\par
\section*{Cyclic \(p\)-groups and finite diagnostics}
For \(G=C_{p^n}\), Notation 8.1 and Proposition 8.3 (pp. 31–32) label \(N_j\) by its index \(p^j\), so \(N_0=G\) and \(N_n=1\). The quotient \(G/N_0\) contributes one cohomological point and each nontrivial cyclic quotient contributes two. The point count is therefore \(2n+1\); when \(n=0\), the space is a singleton.\par
For \(1\le i\le n\), put \(M_i=k(G/N_{i-1})\). Its \(N_j\)-fixed set is empty for \(j<i-1\), and is the whole coset set for \(j\ge i-1\). In the latter case the resulting \(k(C_{p^j})\)-module has the form \[k[t]/(t^{p^{i-1}})\quad\text{over}\quad k[t]/(t^{p^j}).\] It is regular and perfect at \(j=i-1\). At \(j>i-1\), its nonzero dimension is smaller than \(p^j\), so it cannot be free over the local group algebra. Self-injectivity rules out finite projective dimension for this nonprojective module. This gives precisely the support formula used on p. 32: \[\operatorname{supp}(M_i)=\{m_j:j\ge i-1\}\ \cup\ \{\mathfrak p_j:j\ge i\}.\] Combining it with the subgroup-order restriction gives \(\overline{\{\mathfrak p_i\}}\subseteq\{m_{i-1},\mathfrak p_i,m_i\}\). The within-stratum specialization to \(m_i\) is already known. Omitting the remaining edge to \(m_{i-1}\) would disconnect the finite specialization graph. The cited rigid-category splitting principle would then give a nontrivial idempotent in \(\operatorname{End}(\mathbf1)=k\), impossible for a field. This verifies the boundary and connectivity argument conditional on the cited splitting principle and cyclic cohomology input.\par
The native helper’s independently written exact script checked the following finite instances; these corroborate the elementary calculations and do not prove the universal statements.\par
Check | Actual domain | Cases\par
Lemma 5.3 transfer index | \(S_3,\ p=3\); all valid subgroup data | 48\par
Lemma 5.3 transfer index | \(A_4,\ p=2\); all valid subgroup data | 192\par
Cyclic fixed cosets | \(p=2,3,5,\ n\le3\); nontrivial module cases | 60\par
The corrected diagnostic validates subgroup multiplication and inverse closure before applying the hypotheses, and its recorded execution exited with code zero. An earlier reviewer-written diagnostic omitted self-products in subgroup generation and admitted invalid sets. That failed script, result, and explanation are preserved and excluded from paper evidence. The paper is not assigned an error arising from that diagnostic.\par
\section*{Assessment and remaining coverage}
The supported assessment is No Issue Found within the declared scope. Four inspected chains retain their needed hypotheses, and no admissible counterexample or invalid implication was established. No correction is proposed. The finite tests, source-page checks, and exact arguments support this bounded result, not a statement that every theorem in the paper is correct.\par
The full conservativity proof in Section 6, stratification in Section 9, and twisted-cohomology constructions, presentations, closure calculations, and larger examples in Sections 12–18 were not fully verified. External references were not independently audited. No journal-final comparison, later-version/erratum search, author-code reproduction, or author adoption is asserted. The result is an initial AI technical review with native bounded assistance; it is not independent acceptance or a human decision. Scientific decision remains pending.\par
\section*{Supporting checks}
`312ccd6c4a449498a4caa8f2a20567ff4d3057d56a205a810b296d2ac320752b`.\par
The full acquired PDF, source manifest, static author-source tree, fifteen inspected source-page images, detailed proof checks, exact diagnostic scripts/results and revision history are preserved. Availability of author code is not used as a mathematical quality criterion.\par
9 P. Balmer and M. Gallauer, The geometry of permutation modules, arXiv:2307.04398v2, version stamp 21 July 2025, 76 pp. Exact acquired author manuscript; journal-final equivalence unverified.\par
\end{document}
